% CHECKPOINT34: new isolated source; all prior sources unchanged.
\statement{A.106}{Proved}{The actual hexagon in the central congruence quotient and its exact coordinates}
Retain the actual negative full-\(\widehat{GT}\) involution,
canonical \(u\), and marked \(k,d,H_d,q,C,N\) of A.103--A.104.
Thus \(F=\widehat F(x,y)\), \(d=1+2k\), \(k\in2\widehat{\mathbb Z}\), and
\[
 H_d=u(x,y)x^k u(z,x)z^k u(y,z)y^k,\qquad z=(xy)^{-1}.
 \tag{A.106.1}
\]
Let \(C_{(2)}\) be the actual closed subgroup generated by \(X_2,Y_2\)
in \({\rm SL}_2(\mathbb Z_2)\). For the maximal pro-two quotient
\(\pi_2:F\to F^{(2)}\), put
\[
 \begin{gathered}
 N_1=\overline{[N,F]},\qquad
 M_2=\ker(F^{(2)}\to C_{(2)}),\\
 M_{2,1}=\overline{[M_2,F^{(2)}]}.
 \end{gathered}
 \tag{A.106.2}
\]
Then
\[
 \boxed{H_d\in N_1},\qquad
 \boxed{N_1=N\cap F'\cap\pi_2^{-1}(M_{2,1})}.
 \tag{A.106.3}
\]
Consequently the actual compact group \(F/N_1\) embeds continuously,
with its subspace topology, in
\[
 C\times F_{\rm ab}\times(F^{(2)}/M_{2,1}).
 \tag{A.106.4}
\]
Its image is not asserted to be the whole product. No surjectivity
of \(N\to M_2\), or triviality of \(H_d\) in \(F\), is inferred.

\begin{proof}
\emph{The actual congruence image.}
For each prime \(p\), let \(e_p\in\widehat{\mathbb Z}\) be one in
its \(p\)-coordinate and zero in every other coordinate. The elements
\[
 X^{e_p}=I+2e_pE_{12},\qquad Y^{e_p}=I-2e_pE_{21}
\]
belong to \(C\) and are supported at that prime. For odd \(p\),
their powers give every upper and lower elementary matrix over
\(\mathbb Z_p\). These generate \({\rm SL}_2(\mathbb Z_p)\):
a determinant-one matrix has a unit in its first column, so
elementary row operations, with a row exchange if needed, reduce it
to a diagonal matrix. For a unit \(a\), the latter are elementary
products because
\[
 \begin{gathered}
 w(a)=E_{12}(a)E_{21}(-a^{-1})E_{12}(a)
      =\begin{pmatrix}0&a\\-a^{-1}&0\end{pmatrix},\\
 w(a)w(-1)=\operatorname{diag}(a,a^{-1}).
 \end{gathered}
\]
At two the supported generators retain exactly \(C_{(2)}\).
Finite products of these supported factors are dense, and \(C\) is
closed. Hence
\[
 C=C_{(2)}\times C_{\rm odd},\qquad
 C_{\rm odd}=\prod_{p\ {\rm odd}}G_p,\quad
 G_p={\rm SL}_2(\mathbb Z_p).
 \tag{A.106.5}
\]
The group \(C_{(2)}\) is pro-two, since its generators are identity
modulo two. It is not replaced by all of
\({\rm SL}_2(\mathbb Z_2)\). Each \((G_p)_{\rm ab}\) is pro-\(p\):
its finite abelian quotients are generated by images of the two
pro-\(p\) unipotent cyclic subgroups.

\emph{Continuous degree-two cohomology.}
All coefficients in this paragraph are the discrete trivial module
\(\mathcal A=\mathbb Q/\mathbb Z\). The precise new primary input is
\cite[Proposition~5.7(b), printed p.~41]{A_Beyl1986}:
\[
 H^2_{\rm cont}(G_p,\mathcal A)=0.
 \tag{A.106.6}
\]
Only its odd-prime cases are used. Its complete statement and proof
were read in the original GDZ OCR XML. The proof takes the direct
limit of finite-level cohomology and uses vanishing transition maps;
it does not require every finite-level multiplier to vanish.

Here is the product argument at the needed scope. If profinite
\(P,Q\) satisfy \(H^2(Q,\mathcal A)=0\), and every continuous
bicharacter \(P_{\rm ab}\times Q_{\rm ab}\to\mathcal A\) is zero,
then projection induces an isomorphism
\[
 H^2(P,\mathcal A)\xrightarrow{\ \mathrm{inf}\ }
                  H^2(P\times Q,\mathcal A).
 \tag{A.106.7}
\]
Restriction to \(P\) is a left inverse. For surjectivity, subtract the
inflation of the restriction of a given class to \(P\). Its
restrictions to both factors are now trivial. In its continuous
central extension by \(\mathcal A\), choose homomorphic continuous
sections of the two factors. Their commutators form the stated
bicharacter. It vanishes, so the sections commute and their product
splits the extension.

For pro-\(p\) and pro-\(q\) abelian groups with \(p\ne q\), such a
bicharacter vanishes: its image is finite by compactness, and each
value has both \(p\)-power and \(q\)-power order. Thus (A.106.6)--(A.106.7)
give \(H^2=0\) for every finite product of odd factors.
A continuous cocycle on the infinite product has finite image and
depends on finitely many prime factors. Indeed compactness gives
a common open normal coordinate neighborhood on which it is constant;
this neighborhood contains all factors outside a finite set. The
descended function remains a cocycle. Therefore
\(H^2(C_{\rm odd},\mathcal A)=0\).
The same finite-dependence argument makes every bicharacter between
\(C_{(2),\rm ab}\) and \(C_{{\rm odd},\rm ab}\) zero: its order is
both two-primary and odd. It follows that
\[
 H^2(C_{(2)},\mathcal A)
       \xrightarrow{\ \mathrm{inf}\ }H^2(C,\mathcal A)
 \quad\hbox{is an isomorphism}.
 \tag{A.106.8}
\]
In particular no assertion \(H^2(C_{(2)},\mathcal A)=0\) is used.

\emph{Comparison of the actual kernels.}
There is the actual commutative extension diagram
\[
\begin{array}{ccccccccc}
1&\to&N&\to&F&\to&C&\to&1\\
 &&\downarrow&&\downarrow\pi_2&&\downarrow\\
1&\to&M_2&\to&F^{(2)}&\to&C_{(2)}&\to&1 .
\end{array}
 \tag{A.106.9}
\]
Both free groups have continuous \(H^2(-,\mathcal A)=0\).
For the free pro-two group this uses cohomological dimension one
on two-primary torsion and vanishing on odd-primary torsion;
continuous cochains commute with the filtered union of finite
torsion submodules of \(\mathcal A\).
The continuous five-term sequences therefore give surjections
\[
 \begin{gathered}
 H^1(N,\mathcal A)^F\xrightarrow{\operatorname{trg}}
                      H^2(C,\mathcal A),\\
 H^1(M_2,\mathcal A)^{F^{(2)}}
       \xrightarrow{\operatorname{trg}_2}H^2(C_{(2)},\mathcal A),
 \end{gathered}
\]
whose kernels consist of restrictions of characters of \(F\) and
\(F^{(2)}\), respectively. Use the same transgression convention
in both rows. By (A.106.8), for every invariant continuous
\(\chi:N\to\mathcal A\), choose an invariant
\(\chi_2:M_2\to\mathcal A\) with the corresponding transgression.
Naturality under (A.106.9) and exactness give a character
\(\lambda:F\to\mathcal A\) with
\[
 \chi=\chi_2\circ\pi_2|_N+\lambda|_N.
 \tag{A.106.10}
\]
This does not require \(N\to M_2\) to be onto.

\emph{The exact coordinate criterion and the hexagon.}
Commutator functoriality gives
\(N_1\subseteq N\cap F'\cap\pi_2^{-1}(M_{2,1})\).
For the reverse inclusion, take \(h\) in that intersection.
Every summand on the right of (A.106.10) kills \(h\):
\(\chi_2\) kills \(M_{2,1}\) and \(\lambda\) kills \(F'\).
Invariant characters of \(N\) are precisely the characters of
the profinite abelian group \(N/N_1\). They separate points by
finite abelian quotients and their characters into \(\mathcal A\).
Thus \(h\in N_1\), proving the second equality in (A.106.3).
The map (A.106.4) has exactly this kernel; compactness and the
Hausdorff property give the asserted topological embedding.

Finally A.103 gives \(H_d\in N\). The substituted \(u\)-words
lie in \(F'\), and the exponent vector of \(x^kz^ky^k\) is zero,
so \(H_d\in F'\). Crucially A.89, proof part~3, gives the
\emph{whole} image \(\pi_2(H_d)=1\), not only a dyadic
coefficient equality: A.42 applies to this same canonical \(u\),
A.45 identifies its local roots as \(\pm d_2\), and the
prescribed pro-two positive hexagon with root \(d_2\) is exactly
the image of (A.106.1). The coordinate criterion therefore proves
\(H_d\in N_1\). This last input has full-\(\widehat{GT}\) scope;
it is not supplied for an arbitrary negative
\(\widehat{GT}_0\)-involution.
\end{proof}

The theorem retains the actual signed marking and the same canonical
\(u,d\). It neither proves the full hexagon nor identifies any
abstract completion of a discrete congruence kernel.

